\section{Experiments}
\label{sec:experiments}

\subsection{Experimental Setup}
\label{sec:setup}

\paragraph{Models.}
We probe two instruction-tuned models with comparable parameter counts but distinct architectures: Llama-3.1-8B-Instruct (32 layers, hidden dimension 4096) \citep{grattafiori2024llama} and Qwen3-8B (36 layers, hidden dimension 4096) \citep{yang2024qwen25}.

\paragraph{Dataset.}
We use ConflictQA \citep{xie2024adaptive}, a benchmark of context-parametric knowledge conflicts where a provided context contradicts the model's parametric knowledge.
After verifying that each model can correctly answer the parametric-knowledge question without context, we retain 1{,}368 samples for Llama (1{,}096 train / 272 test) and 966 samples for Qwen (773 train / 193 test), split 80/20 with stratification.

\paragraph{Translation.}
We translate all English contexts and questions into Chinese using a Qwen3-8B translation pipeline (EN$\to$ZH), with quality validated by COMET \citep{rei2020comet} back-translation scores.

\paragraph{Probing.}
For each layer $l$, we train a logistic regression probe on the last-token residual stream representation $\mathbf{h}_l$, with binary labels indicating whether the model chose the contextual or parametric answer.
We report balanced accuracy to account for class imbalance.
All confidence intervals (CIs) are computed via bootstrap resampling ($n = 2{,}000$) with paired differences for asymmetry estimation.

\paragraph{Non-conflict baseline.}
To test whether cross-lingual transfer patterns are specific to knowledge conflicts, we run the same probing pipeline on XNLI \citep{conneau2018xnli}, classifying entailment versus contradiction.

\subsection{Within-Language Conflict Detection}
\label{sec:within_lang}

Table~\ref{tab:main_results} reports balanced accuracy for all probe transfer directions.
Both models encode linearly detectable conflict signals in their residual streams.
Llama achieves 77.6\% on EN$\to$EN (peak at layer 28, 88\% depth) and 72.1\% on ZH$\to$ZH (peak at layer 13, 41\% depth).
Qwen reaches 75.3\% on EN$\to$EN (peak at layer 22, 61\% depth) and 72.6\% on ZH$\to$ZH (peak at layer 14, 39\% depth).
All within-language accuracies are significantly above chance (binomial $p < 0.001$).
Figure~\ref{fig:layer_probe} shows the layer-wise probe accuracy curves for both models.
These results confirm that knowledge conflict signals are linearly separable in the residual stream across both architectures and both languages.

\subsection{Cross-Lingual Transfer and Asymmetry}
\label{sec:cross_lingual}

All four cross-lingual transfer directions achieve balanced accuracy above chance (binomial $p < 0.001$; Table~\ref{tab:main_results}), confirming that conflict representations partially overlap across languages.
Both models show a consistent asymmetry: probes trained on Chinese representations transfer to English (ZH$\to$EN) more accurately than the reverse direction (EN$\to$ZH).
For Llama, the asymmetry $\Delta_{\text{asym}} = +7.2$ percentage points (pp), with a 95\% bootstrap CI of [1.1, 13.6] that excludes zero.
For Qwen, $\Delta_{\text{asym}} = +7.2$ pp with CI [$-$1.2, 13.4], which contains zero and is therefore not statistically significant at $\alpha = 0.05$.
The consistent direction of asymmetry across architectures suggests a shared structural property rather than a model-specific artifact (Figure~\ref{fig:asymmetry_forest}).

%% Table 1: Main cross-lingual probe results
\begin{table}[t]
\centering
\caption{Cross-lingual probe balanced accuracy (\%) with bootstrap 95\% CIs ($n = 2{,}000$). $\Delta_{\text{asym}}$ = acc(ZH$\to$EN) $-$ acc(EN$\to$ZH), computed as the mean of paired bootstrap differences. Bold: best per-task direction.}
\label{tab:main_results}
\begin{tabular}{lcc}
\toprule
Direction & Llama-3.1-8B & Qwen3-8B \\
\midrule
\multicolumn{3}{l}{\textit{Conflict Task}} \\
EN $\to$ EN & \textbf{77.6} & \textbf{75.3} \\
ZH $\to$ ZH & 72.1 & 72.6 \\
EN $\to$ ZH & 59.6 & 67.3 \\
ZH $\to$ EN & 71.3 & 73.2 \\
$\Delta_{\text{asym}}$ & +7.2 [1.1, 13.6] & +7.2 [$-$1.2, 13.4] \\
\midrule
\multicolumn{3}{l}{\textit{XNLI Baseline}} \\
EN $\to$ EN & 94.0 & 96.6 \\
ZH $\to$ ZH & 87.5 & 89.5 \\
EN $\to$ ZH & 89.0 & 90.5 \\
ZH $\to$ EN & \textbf{94.9} & \textbf{97.2} \\
$\Delta_{\text{asym}}$ & +6.0 [0.9, 11.3] & +6.7 [2.4, 11.7] \\
\bottomrule
\end{tabular}
\end{table}

\subsection{Non-Conflict Baseline Comparison}
\label{sec:xnli_baseline}

If the ZH$\to$EN transfer advantage were specific to knowledge conflicts, it should be absent or substantially reduced in a non-conflict semantic task.
We test this prediction using XNLI entailment-versus-contradiction probes.
Both models achieve high within-language accuracy on XNLI (Llama 94.0\% EN, 87.5\% ZH; Qwen 96.6\% EN, 89.5\% ZH), confirming that the probing pipeline captures semantic distinctions beyond knowledge conflicts.
The XNLI probes exhibit the same directional asymmetry: Llama $\Delta_{\text{asym}} = +6.0$ pp, CI [0.9, 11.3]; Qwen $\Delta_{\text{asym}} = +6.7$ pp, CI [2.4, 11.7].
Both XNLI CIs exclude zero, and both overlap substantially with the corresponding conflict-task CIs (Figure~\ref{fig:asymmetry_forest}).
The transfer asymmetry is therefore not specific to knowledge conflict representations but reflects a general cross-lingual property of these models: Chinese-trained probes transfer to English more readily than the reverse, regardless of the downstream task.

\subsection{Confound Exclusion}
\label{sec:confounds}

Two potential confounds could inflate cross-lingual probe accuracy: translation quality and lexical overlap between English and Chinese inputs.

\paragraph{Translation quality.}
We stratify test samples by COMET score into terciles (high, medium, low) and compare EN$\to$ZH probe accuracy across strata.
If translation quality drove probe transfer, high-COMET samples would yield higher accuracy.
Table~\ref{tab:confound} shows no such pattern: Qwen exhibits a reversed direction ($\Delta = -5.6$ pp, Fisher $p = 0.695$), and Llama shows a negligible effect ($\Delta = +1.7$ pp, Fisher $p = 0.761$).
Translation quality does not confound the cross-lingual transfer results.

\paragraph{Token overlap.}
Per-sample Jaccard similarity between English and Chinese tokenized inputs averages 3--4\%, consistent with the minimal lexical overlap between the two languages.
Neither model shows a positive correlation between token overlap and probe accuracy: Qwen yields no significant correlation ($p > 0.4$), and Llama shows a weak negative correlation ($r = -0.15$, $p = 0.012$), where higher overlap is associated with \emph{lower} accuracy.
Shared tokens do not drive cross-lingual transfer.

%% Table 2: Confound exclusion
\begin{table}[t]
\centering
\caption{Confound exclusion analysis. COMET tercile: $\Delta$ = balanced accuracy difference between high and low COMET score terciles for EN$\to$ZH probes (Fisher's exact test). Token overlap: Pearson correlation between per-sample Jaccard similarity and EN$\to$ZH probe accuracy.}
\label{tab:confound}
\begin{tabular}{llrl}
\toprule
Confound & Model & Result & $p$ \\
\midrule
COMET tercile & Qwen & $\Delta = -5.6$ pp & 0.695 \\
COMET tercile & Llama & $\Delta = +1.7$ pp & 0.761 \\
\midrule
Token overlap & Qwen & n.s. & $>$0.4 \\
Token overlap & Llama & $r = -0.15$ & 0.012 \\
\bottomrule
\end{tabular}
\end{table}

\subsection{Causal Evidence via Representation Steering}
\label{sec:steering}

The probing results establish that conflict signals are detectable, but detection does not imply causal relevance.
We test whether the probe weight direction $\mathbf{w}_{\text{probe}}$ causally influences conflict resolution by injecting $\alpha \cdot \mathbf{w}_{\text{probe}}$ into the last-token residual stream at a target layer during inference.
We measure the flip rate (fraction of samples that switch answer category) and compare against a random-direction control of equal magnitude via Fisher's exact test.

\paragraph{L13 (mid-layer, 41\% depth).}
Steering at layer 13 produces significantly more behavioral change than random controls in both directions (Table~\ref{tab:steering}).
Positive steering ($\alpha = +5$) yields a flip rate of 19.9\% versus 6.6\% for the random control (Fisher $p = 0.002$), with all 4 directional flips moving from parametric to contextual answers.
Negative steering ($\alpha = -20$) yields 36.0\% versus 17.6\% (Fisher $p < 0.001$), with 5 of 6 directional flips moving from contextual to parametric.
Both results survive Bonferroni correction for multiple comparisons.
The bidirectionality confirms that the probe weight vector encodes a meaningful axis of conflict resolution: positive perturbations push toward contextual answers, and negative perturbations push toward parametric answers.
Steering at nearby L16 (50\% depth, just three layers away) showed no significant effect (all $p > 0.12$), indicating a remarkably narrow causal window.

\paragraph{Negative layers.}
Steering at L28 (the probing peak layer, 88\% depth) produces no significant effect with either last-token or all-position injection (all Fisher $p > 0.3$; Table~\ref{tab:steering}).
L18 (the MLP probe peak, 56\% depth) also shows no causal influence (all Fisher $p > 0.44$).

\paragraph{$\alpha = +20$ reversal.}
At extreme perturbation magnitude ($\alpha = +20$), the control condition produces more flips than the steering condition ($p = 0.053$).
This indicates coherence collapse at large perturbation magnitudes, constraining the effective steering range.

\paragraph{Detection $\neq$ Decision.}
The detectable conflict signal peaks at late layers (L28, 88\% depth), but causal influence on conflict resolution is localized to the mid-layer (L13, 41\% depth).
Probing accuracy and causal relevance thus identify different computational stages: late layers encode the outcome of conflict resolution in a linearly readable form, while mid-layers host the computation that determines the outcome.

%% Table 3: Steering results
\begin{table}[t]
\centering
\caption{Representation steering results (Llama-3.1-8B). Flip rate: fraction of test samples switching answer category. Control: random-direction perturbation of equal magnitude. Significant results ($p < 0.05$) in bold.$^{\dagger}$}
\label{tab:steering}
\small
\begin{tabular}{lcrrrrl}
\toprule
Layer & Depth & $\alpha$ & Steer (\%) & Ctrl (\%) & Fisher $p$ & Dir. \\
\midrule
\textbf{L13} & \textbf{41\%} & $\mathbf{+5}$ & \textbf{19.9} & \textbf{6.6} & $\mathbf{0.002}$ & \textbf{4 p$\to$c} \\
\textbf{L13} & \textbf{41\%} & $\mathbf{-20}$ & \textbf{36.0} & \textbf{17.6} & $\mathbf{<0.001}$ & \textbf{5 c$\to$p} \\
L16 & 50\% & $+5$ & 8.8 & 8.1 & 1.000 & \\
L16 & 50\% & $+10$ & 18.4 & 11.0 & 0.123 & \\
L16 & 50\% & $+20$ & 22.8 & 24.3 & 0.886 & \\
L18 & 56\% & $+2$ & 3.7 & 1.5 & 0.447 & \\
L28 & 88\% & $-20$ & 14.7 & 16.9 & 0.740 & \\
L28 (all) & 88\% & $-20$ & 11.8 & 11.8 & 1.000 & \\
\bottomrule
\end{tabular}
\begin{flushleft}
\footnotesize $^{\dagger}$ At $\alpha = +20$, the control flip rate exceeds steering ($p = 0.053$), indicating coherence collapse at extreme perturbation magnitudes.
\end{flushleft}
\end{table}
